\section{Implementation}
\label{sec:implementation}

The \ac{bpfima} framework leverages \ac{eBPF} technology to extend Linux kernel integrity measurement capabilities.
It introduces container-aware monitoring and dynamic policy enforcement.
System composition includes a custom kernel module, \ac{eBPF} programs hooked into \ac{LSM} points, and a user space control utility.

\subsection{Kernel Module}
Core functionality resides in a Linux kernel module establishing the integrity measurement infrastructure.
Upon initialization, it creates a \texttt{SecurityFS} interface at \texttt{/sys/kernel/security/bpfima/}.
This interface serves as the primary communication channel between kernel and user space.
It exposes global status, policy configurations, and a hierarchical view of measurements organized by container namespaces.

Central to the design is a non-binary Merkle tree.
Unlike traditional flat lists in \ac{IMA}, \ac{bpfima} assigns a dedicated leaf node to each active container.
Each leaf corresponds to a unique namespace identifier.
This structure isolates measurement chains.
Therefore, the integrity state of one container does not obscure others.
The Merkle root hash acts as a global system state.
It is periodically extended into the \ac{TPM} \ac{PCR} 23 to provide a hardware root of trust.
This behavior is configurable through policies.

Custom \ac{BTF}-enabled \textit{kfuncs} are registered by the module.
These functions expose internal kernel operations to \ac{eBPF} programs.
Such design bypasses standard \ac{eBPF} helper limitations.
It enables direct interaction with the \ac{IMA} subsystem and \ac{TPM} drivers from \ac{LSM} hooks.
Key functions include \texttt{bpfima\_measurement\_extend} for recording events and container handling.

\subsection{User Space Tool}
\texttt{bpfima-tool} is the user space control component.
Developed in C, it uses \texttt{libbpf} for robust \ac{CO-RE} (Compile Once - Run Everywhere) support.
Primary responsibilities include loading \ac{eBPF} bytecode into the kernel and attaching programs to hooks.
It also manages the entire monitoring process lifecycle, measuring loaded/unloaded \ac{eBPF} programs.

Two operational modes are available.
Management tasks like loading policies or checking status use an interactive command-line interface.
Continuous monitoring relies on a daemon mode.
Bpfima-tool facilitates through bpf maps the storage of lists at user side.
The tool handles signal processing (\texttt{SIGINT}, \texttt{SIGTERM}) for graceful shutdown.
Cleanup routines ensure that BPF maps are unpinned and resources are released upon exit.

\subsection{Policies}
A dynamic, granular policy engine decouples security logic from enforcement mechanisms.
Policies are defined in human-readable YAML format.
The user space tool translates these into \ac{eBPF} Maps, specifically \texttt{BPF\_MAP\_TYPE\_ARRAY} for efficient O(1) lookups.
Kernel-side hooks read these maps directly.
This architecture allows administrators to update security policies at runtime.
No kernel module reload or application restart is required.

Three distinct layers structure the policy model.
Global Policy configuration controls system-wide behavior.
This includes the master enable switch, minimum file size thresholds, and logging verbosity.
Underneath, a filtering layer manages exclusion patterns.
It acts as a deny-list for control groups (cgroups) and file paths.
Noise from infrastructure processes or ephemeral filesystems is thus filtered out.
Finally, the Hook Configuration layer provides fine-grained control.
Individual \ac{LSM} hooks can be independently enabled, disabled, or set to "track-only" mode.

Integrity of the policy itself is strictly monitored.
Any active policy map modification is treated as a measureable event.
New configurations are serialized, hashed, and extended into the global Merkle tree root.
Consequently, the security configuration state is cryptographically bound to the hardware root of trust.

A typical policy configuration enables comprehensive monitoring while reducing noise.
For instance, a standard deployment might enforce:
\begin{verbatim}
policy:
  enabled: true
  log_level: 2
filters:
  proc_sys: false    # Monitor /proc and /sys
  tmp_files: false   # Monitor /tmp
actions:
  extend_tpm: true
  track_container: true
hooks:
  lsm_bprm_check_security:
    enabled: true
\end{verbatim}
This example effectively activates the system with informational logging.
By disabling the `proc_sys` and `tmp_files` filters, it mandates coverage of all file operations, including those in temporary directories often targeted by attackers.
The actions ensure that every event is both recorded for software audit and cryptographically extended into the \ac{TPM}, satisfying strict integrity requirements.
Reference hooks like \texttt{lsm\_bprm\_check\_security} are explicitly activated to intercept binary execution, ensuring that no code runs inside a container without being measured.

\subsection{eBPF Functions and Library}
\texttt{headers\_bpf.h} is a specialized library provided to facilitate secure integrity extensions.
It abstracts the complexity of raw \ac{eBPF} interactions.
Standard data structures for policies and measurements are defined here.
Inline helper functions like \texttt{bpfima\_should\_process} standardizes "policy check" logic across all hooks.
This ensures consistent filtering rule application.

Core \textit{kfuncs} exported by the kernel module back the library.
These perform privileged operations typically unavailable to unprivileged \ac{eBPF} bytecode:

\textbf{Measurement Extension:} \texttt{bpfima\_measurement\_extend} is the primary interface for recording integrity violations.
When an \ac{eBPF} program detects a relevant event, it invokes this function.
Arguments include the event name, target namespace ID, and dependency metadata.
Atomicity is handled entirely within the kernel.
The function locks the container's leaf node, calculates the new hash chain value, and updates the global Merkle root.
Optional hardware TPM extension is also triggered here.

\textbf{File Hashing:} Standard \ac{eBPF} programs encounter limits when computing cryptographic hashes of large files.
The \texttt{bpfima\_file\_hash} kfunc solves this.
It accepts a scalar representation of a `struct file` pointer.
This pointer is safely cast to avoid verifier rejection.
Actual SHA-256 computation is then delegated to the kernel's native \ac{IMA} subsystem.

\textbf{Trust Anchor Retrieval:} Conditional security logic is enabled by functions like \texttt{bpfima\_tpm\_get\_pcr\_value}.
\ac{eBPF} programs can read current PCR register values directly from TPM hardware.
Developers can thus write "smart" hooks.
For example, access to sensitive keys might be permitted only if the system's integrity state matches a known-good policy.
